% ============================================================
%  Section 6: Training and Inference
%  H&E to Sirius Red Virtual Staining Paper
% ============================================================

%\subsection{Training and Inference}
%\label{sec:training_inference}

\textbf{Training Protocol.}
All models are trained with the Adam optimiser~\cite{kingma2015adam} (full hyperparameters
in \suppsec{3}). We count training in optimiser steps rather than epochs,
ensuring fair comparison across data-fraction configurations with different tile counts.
CycleGAN, UNIT, MUNIT, DCLGAN, and CycleDiffusion are each trained for $750{,}000$~steps.
UVCGAN follows a two-stage schedule: $250{,}000$~steps of masked-image pretraining
followed by $500{,}000$~steps of cycle-consistent finetuning initialised from the
Stage~1 checkpoint.

\noindent\textbf{Scaling Study Design.}
We conduct a full factorial experiment over the six architectures of
Section~\ref{sec:models}, three generator capacity scales, and three training data
fractions.

\noindent\textbf{Generator capacity.}
Each model is trained at three generator sizes: Small~(S, ${\sim}10$~million parameters),
Medium~(M, ${\sim}50$~million), and Large~(L, ${\sim}100$~million), by adjusting the
base channel width and depth (or equivalent per model family; see \suppsec{2}).
All size variants within a family share the same training protocol, data, and random
seed, so any difference reflects capacity alone.

\noindent\textbf{Training data fraction.}
The training set is partitioned into three nested subsets comprising 25\%, 50\%, and
100\% of the 30 training paired specimens (60 WSIs), constructed via contiguous WSI ranges
to avoid data leakage.
Smaller subsets are strict subsets of larger ones, so differences across data fractions
are not confounded by variation in which slides are included.
The same fixed test set of $N{=}5$ held-out paired specimens (10 WSIs) is used for
evaluation across all configurations.
Together, these axes yield $6\text{ models} \times 3\text{ sizes} \times 3\text{ data
fractions} = \mathbf{54}$ experimental configurations, each trained and evaluated
independently.

\noindent\textbf{Inference.}
At test time, all GAN-based models perform a single forward pass through the trained
$G_{A \to B}$ generator per tile.
CycleDiffusion requires two sequential DDIM passes per tile~\cite{song2021ddim}: the
source-domain DDPM encodes the source tile into the shared noise latent via DDIM
inversion over $T_{\text{encode}}{=}$200~steps, and the target-domain DDPM then decodes
it into the target domain via DDIM sampling over $T_{\text{sample}}{=}$200~steps.

% ============================================================
%  New bibliography entries for this section
%  (merge with references.bib)
% ============================================================
%
% @inproceedings{kingma2015adam,
%   author    = {Kingma, Diederik P. and Ba, Jimmy},
%   title     = {Adam: A Method for Stochastic Optimization},
%   booktitle = {ICLR},
%   year      = {2015},
%   eprint    = {1412.6980}
% }
%
% (Shared entries already listed:
%  song2021ddim)
